\appendix
\section*{Artifact Appendix}
\addcontentsline{toc}{section}{Artifact Appendix}

\subsection*{A.1 Artifact Overview}
This appendix describes how to reproduce the project. The artifact contains the Spring Boot consumer, RabbitMQ and Kubernetes manifests, HPA and KEDA scaling policies, Python load-generation scripts, shell scripts for deployment and data collection, raw experimental CSV files, and Python analysis scripts that regenerate the figures and summary table.

\subsection*{A.2 Hardware and Software Dependencies}
\begin{itemize}
  \item \textbf{Cloud host}: AWS EC2 \texttt{m5.large} with Ubuntu~22.04~LTS. A \texttt{t3.medium} can be used for smoke testing, but the reported experiments use \texttt{m5.large}.
  \item \textbf{Cluster stack}: K3s v1.30, Helm~3, \texttt{metrics-server} v0.7, KEDA v2.13, and RabbitMQ \texttt{3.12-management}.
  \item \textbf{Build tools}: Java~17, Maven~3.9, Docker~24.x or compatible Docker Buildx environment.
  \item \textbf{Load-test tools}: Python~3.10+ with \texttt{pika}, \texttt{click}, and \texttt{PyYAML} as listed in the load-test requirements file.
  \item \textbf{Analysis tools}: Python~3.10+ with \texttt{pandas} and \texttt{matplotlib} as listed in the analysis requirements file.
\end{itemize}

\subsection*{A.3 Repository}
The public repository is:
\begin{center}
\url{https://github.com/QinHaoting/CS5296-Group5-Autoscaling}
\end{center}

Key directories are:
\begin{itemize}
  \item \texttt{consumer/}: Spring Boot RabbitMQ consumer.
  \item \texttt{k8s/}: namespaces, RabbitMQ, HPA baseline, and KEDA manifests.
  \item \texttt{load-test/}: workload producer and YAML traffic patterns.
  \item \texttt{scripts/}: setup, deployment, collection, run, and teardown scripts.
  \item \texttt{results/}: raw CSV data, generated figures, and \texttt{summary.csv}.
  \item \texttt{analysis/}: scripts for metric extraction and plotting.
\end{itemize}

\subsection*{A.4 Build and Deployment}
\begin{enumerate}
  \item Build and push the consumer image:
\begin{verbatim}
docker buildx build --platform linux/amd64 \
  -t <dockerhub>/cs5296-consumer:v1.0 --push consumer/
\end{verbatim}
  \item On the EC2 node, install the cluster dependencies and base services:
\begin{verbatim}
bash scripts/setup.sh
\end{verbatim}
  \item Deploy both experiment groups:
\begin{verbatim}
export CONSUMER_IMAGE=<dockerhub>/cs5296-consumer:v1.0
bash scripts/deploy-baseline.sh
bash scripts/deploy-keda.sh
\end{verbatim}
  \item Confirm that RabbitMQ, the consumers, HPA, and KEDA are running:
\begin{verbatim}
kubectl get pods -A
kubectl -n baseline get hpa
kubectl -n keda get scaledobject,hpa
\end{verbatim}
\end{enumerate}

\subsection*{A.5 Running the Experiments}
Run one trial for each strategy:
\begin{verbatim}
bash scripts/run-experiment.sh baseline 1
bash scripts/run-experiment.sh keda 1
\end{verbatim}

Run the full three-trial batch:
\begin{verbatim}
for run in 1 2 3; do bash scripts/run-experiment.sh baseline "$run"; done
for run in 1 2 3; do bash scripts/run-experiment.sh keda "$run"; done
\end{verbatim}

Each trial resets the Deployment to one replica, purges the target queue, records a short baseline, runs the burst producer, and keeps collecting metrics during the post-burst drain period. The full six-trial batch takes about 49 minutes, excluding manual setup time.

\subsection*{A.6 Reproducing the Figures}
After the raw CSV files exist, regenerate the summary and figures with:
\begin{verbatim}
python analysis/plot.py
\end{verbatim}
Expected outputs include \texttt{summary.csv}, three per-run timeline figures, and the aggregate figures for responsiveness, scale-down behaviour, and pod-seconds overshoot.

\subsection*{A.7 Expected Results}
The regenerated \texttt{summary.csv} should show that KEDA reacts approximately 10~s faster than the CPU-based HPA, drains the queue approximately 12~s sooner, and triggers scale-down during the trial. The exact numbers may vary slightly with EC2 performance, RabbitMQ startup time, and producer-side timing, but the qualitative result should remain: queue-length-based scaling reacts earlier and releases capacity earlier for this queue-backed workload.

\subsection*{A.8 Troubleshooting}
Additional operational notes are in \texttt{docs/RUNBOOK.md}. Common checks include verifying the RabbitMQ NodePort, confirming that the two queues are empty before each trial, ensuring \texttt{metrics-server} is reporting CPU metrics, and checking that the KEDA \texttt{ScaledObject} has an active RabbitMQ trigger.
